% Segunda lectura — Bloque algebraico-geométrico: productos, determinante y autovectores
\section{Bloque algebraico-geométrico: productos, determinante y autovectores}\label{cap:11}\label{sec:bloque-algebraico}

\textbf{Dictado} (clase del 18 de septiembre). El objetivo declarado por el docente: que la
parte geométrica y la parte algebraica ``se conecten'', y entender por qué el determinante
tiene las propiedades ``mágicas'' que tiene. Lo que sigue reconstruye ese razonamiento en el
orden en que se dio.

\section{Producto interno: incidencia}\label{sec:interno}

\textbf{Desarrollo.} El producto interno sirve para dos cosas: proyectar, y decidir si dos
vectores son ortogonales (\(u\cdot v=0\)). En coordenadas homogéneas aparece un tercer uso.
Un punto se representa por una terna \(p=(x,y,w)^T\) y una recta \(ax+by+c=0\) por otra
terna \(\ell=(a,b,c)^T\). Son objetos distintos (el docente los dibuja con símbolos
distintos), pero ambos son vectores de \(\R^3\), y
\[
  \boxed{\ \ell\cdot p=ax+by+cw=0\quad\Longleftrightarrow\quad p\ \text{está en}\ \ell\ }.
\]
La ecuación cartesiana de la recta es literalmente ``el producto interno con la normal es
cero''. Ni \(\ell\) ni \(p\) cambian de significado si se los multiplica por un número, y la
condición tampoco.

\section{Producto vectorial: unir dos puntos, cortar dos rectas}\label{sec:vectorial}

\textbf{Desarrollo.} Queremos una operación que, dados dos puntos \(p,q\), devuelva la recta
\(\ell=p\times q\) que los une. ¿Qué tiene que cumplir? Que \(p\) esté en esa recta y \(q\)
también:
\[
  (p\times q)\cdot p=0,\qquad (p\times q)\cdot q=0 .
\]
Por lo tanto la operación debe producir un vector \emph{ortogonal a sus dos argumentos}. Y
más: si pongo dos veces el mismo punto no hay recta que determinar, así que \(p\times p=0\).
Estas exigencias, junto con la linealidad en cada argumento, \emph{obligan} a que la
operación sea el producto vectorial de \(\R^3\):
\[
  p\times q=\begin{pmatrix}p_2q_3-p_3q_2\\ p_3q_1-p_1q_3\\ p_1q_2-p_2q_1\end{pmatrix}.
\]

\begin{resultado}[Las dos propiedades que lo definen]
\begin{itemize}
  \item \textbf{Linealidad} en cada argumento: \((2p)\times q=2\,(p\times q)\),
  \((p+p')\times q=p\times q+p'\times q\).
  \item \textbf{Antisimetría}: \(p\times q=-\,q\times p\). Es equivalente a \(p\times p=0\):
  el único vector igual a su opuesto es el nulo.
\end{itemize}
Son exactamente las propiedades del \emph{área orientada} de un paralelogramo (duplicar un
lado duplica el área; cambiar el orden de los lados cambia la orientación; dos lados
paralelos no encierran área). Por eso \(\lVert p\times q\rVert\) es el área del paralelogramo
de lados \(p,q\).
\end{resultado}

Por simetría de la ecuación de incidencia, la misma operación resuelve el problema dual:
\emph{el punto donde se cortan dos rectas es} \(\ell\times m\). Dos rectas paralelas
\((a,b,c)\) y \((a,b,c')\) dan \(\ell\times m=(c'-c)\,(b,-a,0)^T\): un punto ideal, su
dirección común.

\section{Tres puntos alineados: el determinante}\label{sec:determinante}

\textbf{Desarrollo.} ¿Está el punto \(r\) en la recta que une \(p\) y \(q\)? Se forma la
recta y se pregunta por la incidencia:
\[
  (p\times q)\cdot r=0 .
\]
Esta cantidad es simétrica en los tres puntos (si \(r\) está en \(pq\), entonces \(p\) está
en \(qr\) y \(q\) en \(rp\)), se anula si dos de ellos coinciden, es lineal en cada uno y
cambia de signo al intercambiar dos. Es, término a término, el determinante:
\[
  (p\times q)\cdot r=\det\begin{pmatrix}p_1&p_2&p_3\\q_1&q_2&q_3\\r_1&r_2&r_3\end{pmatrix}.
\]

\begin{resultado}[El determinante es la única operación posible]
Una operación de tres vectores de \(\R^3\) que sea lineal en cada argumento y que se anule
cuando dos argumentos coinciden es, a menos de una constante, el determinante. Vale cero
exactamente cuando los tres vectores son linealmente dependientes (coplanares como vectores;
alineados como puntos homogéneos) y su valor absoluto es el volumen del paralelepípedo que
forman. Intercambiar dos filas cambia el signo: es la antisimetría otra vez.
\end{resultado}

\paragraph{Verificación numérica hecha en clase.} Con \(p=(1,2,3)\), \(q=(4,5,-1)\),
\(r=(0,0,1)\):
\[
  p\times q=\begin{pmatrix}2\cdot(-1)-3\cdot5\\ 3\cdot4-1\cdot(-1)\\ 1\cdot5-2\cdot4\end{pmatrix}
  =\begin{pmatrix}-17\\13\\-3\end{pmatrix},\qquad
  (p\times q)\cdot r=-3,
\]
y el determinante de la matriz de filas \(p,q,r\), desarrollado por la tercera fila, es
\(1\cdot(1\cdot5-2\cdot4)=-3\). Coinciden, como debía ser: desarrollar por una fila
\emph{es} hacer un producto vectorial con las otras dos y luego un producto interno.

\section{Producto exterior: la misma idea en cualquier dimensión}\label{sec:exterior}

\textbf{Desarrollo.} El producto vectorial existe sólo en \(\R^3\). La operación que
generaliza la idea a cualquier dimensión es el \emph{producto exterior} de Grassmann,
\(u\wedge v\), definido únicamente por linealidad y antisimetría (\(u\wedge u=0\)). Su
resultado no es un vector sino un \emph{bivector}: un elemento de área orientada.

Escribiendo \(u=ae_1+be_2+ce_3\) y \(v=de_1+ee_2+fe_3\) y distribuyendo, los términos
\(e_i\wedge e_i\) se anulan y \(e_2\wedge e_1=-e_1\wedge e_2\), de modo que
\[
  u\wedge v=(ae-bd)\,e_1\wedge e_2+(af-cd)\,e_1\wedge e_3+(bf-ce)\,e_2\wedge e_3 .
\]
Son los mismos tres números del producto vectorial; sólo cambia el signo del término del
medio, porque la regla de la mano derecha asigna \(e_1\times e_3=-e_2\). Los ``planos
elementales'' \(e_1\wedge e_2\), \(e_1\wedge e_3\), \(e_2\wedge e_3\) son la base de las
áreas orientadas de \(\R^3\).

Operando con un tercer vector \(w=ge_1+he_2+ie_3\) sólo sobreviven los términos que completan
las tres direcciones, y reordenando (cada intercambio de vecinos cambia el signo) todo cae
sobre un único elemento de volumen:
\[
  (u\wedge v)\wedge w=\det[u\ v\ w]\;e_1\wedge e_2\wedge e_3 .
\]
De ahí la regla de los signos \(+,-,+\) del desarrollo por menores.

\begin{idea}[Contar subespacios]
En \(\R^3\) hay \(3\) direcciones de vectores, \(3\) planos elementales y \(1\) volumen; por
eso el resultado final es un número. En \(\R^4\) hay \(4\) direcciones, \(6\) planos, \(4\)
volúmenes tridimensionales y \(1\) hipervolumen. La operación es mecánica en cualquier
dimensión, y eso es lo que la hace útil: permite buscar un vector ortogonal a un subespacio
en \(\R^n\), que es lo que se necesita para autovectores de matrices grandes.
\end{idea}

\paragraph{Ortogonal es lo que maximiza el volumen.} Entre todos los vectores de norma
\(1\), el que forma mayor volumen con un plano dado es el perpendicular al plano; los otros
producen un prisma ``inclinado'' de menor altura. El docente adelantó que de esta idea sale
la regla para encontrar el vector ortogonal a un subespacio en cualquier dimensión.

\section{Núcleo de una matriz y determinante nulo}\label{sec:nucleo}

\textbf{Desarrollo.} El \emph{núcleo} de \(A\) es el conjunto de vectores \(v\) con
\(Av=0\). Si se piensa \(Av\) fila por fila, cada coordenada del resultado es el producto
interno de una fila de \(A\) con \(v\). Que las tres den cero significa que \(v\) es
ortogonal a las tres filas: las tres filas están en el plano de normal \(v\), son linealmente
dependientes, y por lo tanto
\[
  \boxed{\ \text{núcleo}(A)\neq\{0\}\quad\Longleftrightarrow\quad\det A=0\ }.
\]
Es la propiedad del determinante que se usa para encontrar autovalores.

\section{Autovalores y autovectores}\label{sec:autovalores}

\textbf{Dictado.} \(v\neq0\) es \emph{autovector} de \(A\) con \emph{autovalor}
\(\lambda\) si \(Av=\lambda v\): la transformación lo deja sobre su propia recta. En
coordenadas homogéneas eso es exactamente ``la transformación deja fijo el punto que \(v\)
representa''.

\paragraph{El pase de magia.} \(Av=\lambda v\) equivale a \(Av-\lambda Iv=0\), es decir
\[
  (A-\lambda I)\,v=0 :
\]
\(v\) está en el núcleo de \(A-\lambda I\). Para que ese núcleo no sea trivial,
\[
  \boxed{\ \det(A-\lambda I)=0\ }\qquad\text{(polinomio característico)} .
\]
Para \(A\) de \(3\times3\) es un polinomio de grado \(3\) en \(\lambda\).

\begin{resultado}[Cuántos autovalores reales]
\begin{itemize}
  \item Un polinomio real de grado \(3\) tiene tres raíces, todas reales, o una real y dos
  complejas conjugadas. Nunca tres complejas: las raíces complejas de un polinomio real
  vienen de a pares conjugados.
  \item A cada autovalor real le corresponde un autovector real. A un autovalor complejo no
  le corresponde ningún vector de \(\R^3\).
  \item Una matriz de rotación de \(\R^3\) tiene siempre exactamente un autovalor real, que
  vale \(1\): su autovector es el eje. Los otros dos son complejos.
\end{itemize}
\end{resultado}

\begin{idea}[Método del producto cruz para el autovector]
Encontrado \(\lambda\), las filas de \(A-\lambda I\) son todas ortogonales al autovector.
Entonces el autovector es el producto vectorial de dos de esas filas:
\[
  v=(\text{fila}_i)\times(\text{fila}_j) .
\]
Si esas dos filas son paralelas y el producto da cero, se toma otro par. Si \emph{todos} los
pares dan cero, \(\lambda\) es raíz doble y hay todo un plano de autovectores (en
homogéneas: toda una recta de puntos fijos). No hace falta resolver el sistema: es la
misma operación de ``unir dos rectas en su punto común'' de la sección anterior.
\end{idea}

\subsection{Ejemplo 1: reflexión respecto del plano \texorpdfstring{$xy$}{xy}}

\(A=\operatorname{diag}(1,1,-1)\) conserva \(x,y\) e invierte \(z\).
\[
  \det(A-\lambda I)=(1-\lambda)^2(-1-\lambda)=0
  \quad\Longrightarrow\quad\lambda=1\ (\text{doble}),\ \lambda=-1 .
\]
Para \(\lambda=-1\): \(A+I=\operatorname{diag}(2,2,0)\), y el producto cruz de las dos
primeras filas es \((2,0,0)\times(0,2,0)=(0,0,4)\): el eje \(z\), como se esperaba, con
\(A(0,0,4)^T=(0,0,-4)^T\). Para \(\lambda=1\): \(A-I=\operatorname{diag}(0,0,-2)\), y
cualquier par de filas da producto cruz nulo. Es la raíz doble: todo el plano \(xy\) es de
autovectores, porque la reflexión lo deja quieto punto a punto.

\subsection{Ejemplo 2: eje de una rotación compuesta}

Se compone una rotación de \(90^\circ\) alrededor de \(y\) seguida de una de \(90^\circ\)
alrededor de \(z\):
\[
  R=R_z(90^\circ)\,R_y(90^\circ)=
  \begin{pmatrix}0&-1&0\\1&0&0\\0&0&1\end{pmatrix}
  \begin{pmatrix}0&0&1\\0&1&0\\-1&0&0\end{pmatrix}
  =\begin{pmatrix}0&-1&0\\0&0&1\\-1&0&0\end{pmatrix}.
\]
Es una rotación, así que \emph{sin calcular el polinomio} se sabe que \(\lambda=1\) es
autovalor y que su autovector es el eje. Filas de \(R-I\): \((-1,-1,0)\), \((0,-1,1)\),
\((-1,0,-1)\). Producto cruz de las dos primeras:
\[
  (-1,-1,0)\times(0,-1,1)=(-1,\,1,\,1),
  \qquad
  R\begin{pmatrix}-1\\1\\1\end{pmatrix}=\begin{pmatrix}-1\\1\\1\end{pmatrix}\ \checkmark
\]
El eje es la dirección \((-1,1,1)\). Los otros autovalores son \(-\frac12\pm\frac{\sqrt3}2 i\):
un giro de \(120^\circ\) alrededor de ese eje. Con ángulos genéricos (por ejemplo
\(R_z(30^\circ)R_y(50^\circ)\)) la receta es idéntica y se hace en planilla o en NumPy: el
eje resulta \((-0.2263,\,0.8446,\,0.4853)\), y \(Rv=v\) se verifica multiplicando.

\begin{atencion}[Sobre los ``números lindos'']
Los ejercicios de álgebra están diseñados para que den enteros. Una matriz genérica, como
la de una rotación de \(37^\circ\) o la homografía de un pentágono regular, da raíces
irracionales y autovalores con muchos decimales. Eso no es un error: es lo normal. Por eso
el método debe funcionar en planilla o en código, no sólo a mano.
\end{atencion}

\section{Formas lineales: la recta como máquina de medir}\label{sec:formas-lineales}

\textbf{Desarrollo} (a partir de los resúmenes de clase 4 y 9 del profesor). Hasta acá los
vectores fueron flechas: cosas que se suman y se estiran. Hay otro objeto, igual de simple,
que en los libros aparece tarde y que el profesor prefiere presentar apenas definido el
espacio: la \emph{forma lineal}, una función que toma un vector y devuelve un número,
respetando sumas y escalas.

\paragraph{Primero en \(\R\).} Una forma lineal en la recta es multiplicar por un número:
\(\varphi(x)=2x\). Es una regla graduada: los puntos donde \(\varphi\) vale
\(0,1,2,3,\ldots\) están equiespaciados sobre la recta, a distancia \(\tfrac12\). No hay
nada más que saber. Al pasar al plano la regla sigue siendo la misma, pero ahora hay
direcciones que la regla no mide.

\paragraph{Un ejemplo con números chicos.} Sea
\[
  \varphi(x,y)=2x+3y .
\]
Con \(v=(1,0)\) da \(2\); con \((0,1)\) da \(3\); con \((1,1)\) da \(5\); con \((3,-2)\)
da \(0\). Es lineal: \(\varphi(v+w)=\varphi(v)+\varphi(w)\) y
\(\varphi(\lambda v)=\lambda\,\varphi(v)\). Y basta conocer \(\varphi\) en los dos vectores
de la base para conocerla en todos: \(\varphi(x,y)=x\,\varphi(1,0)+y\,\varphi(0,1)\).

\paragraph{Hecho geométrico: una regla graduada tirada en el plano.} Los puntos donde
\(\varphi\) vale lo mismo forman una recta, \(2x+3y=c\), y al variar \(c\) se obtiene una
familia de rectas paralelas equiespaciadas. La forma lineal \emph{es} esa familia: para
medir un punto, se mira qué recta de la familia pasa por él. La recta \(c=0\) pasa por el
origen y contiene la dirección \((3,-2)\), que es lo que la forma ``no ve'' (su núcleo).
El vector \((2,3)\) es perpendicular a todas las rectas de nivel y apunta hacia donde
\(\varphi\) crece.

\begin{center}
\begin{tikzpicture}[scale=.78]
  \begin{scope}
    \clip (-3.2,-3.2) rectangle (4.2,3.2);
    \foreach \c in {-6,-4,-2,0,2,4,6,8}
      \draw[teal,line width=.7pt] (-3.2,{(\c+6.4)/3})--(4.2,{(\c-8.4)/3});
    \draw[teal,line width=1.2pt] (-3.2,{6.4/3})--(4.2,{-8.4/3});
    \draw[axisline] (-3.2,0)--(4.2,0);
    \draw[axisline] (0,-3.2)--(0,3.2);
    \draw[rust,line width=1pt,-{Latex[length=2.4mm]}] (0,0)--(1,1.5)
      node[right,font=\small,labelback] {\((2,3)\)};
    \draw[tealdark,line width=1pt,-{Latex[length=2.4mm]}] (0,0)--(1.5,-1)
      node[below right,font=\small,labelback] {\((3,-2)\): \(\varphi=0\)};
    \fill[ink] (1,1) circle (2pt) node[above left,font=\small,labelback] {\((1,1)\): \(\varphi=5\)};
  \end{scope}
  \foreach \c/\y in {8/-0.13,6/-0.8,4/-1.47,2/-2.13}
    \node[right,font=\scriptsize,text=teal] at (4.2,\y) {\(c=\c\)};
  \foreach \c/\y in {0/2.13,-2/1.47,-4/0.8,-6/0.13}
    \node[left,font=\scriptsize,text=teal] at (-3.2,\y) {\(c=\c\)};
\end{tikzpicture}
\end{center}

\paragraph{Las selectoras y el espacio de las formas.} Las formas más básicas son las que
seleccionan una coordenada: \(\mathsf{x}(x,y)=x\) y \(\mathsf{y}(x,y)=y\). Cualquier otra
es combinación de ellas: \(\varphi=2\,\mathsf{x}+3\,\mathsf{y}\). Las formas lineales se
suman y se multiplican por escalares, y el resultado sigue siendo una forma lineal: forman
un espacio vectorial, el \emph{dual} del plano. En la base canónica una forma se describe
con dos números, igual que un vector, pero conviene escribirla como \emph{fila}:
\[
  \varphi(v)=\begin{pmatrix}2&3\end{pmatrix}\begin{pmatrix}x\\y\end{pmatrix}
  =(2,3)\cdot v .
\]
La fila por la columna es el producto interno con el vector \((2,3)\): la forma se
``disfraza'' de vector gracias al producto interno. En coordenadas homogéneas esto ya se
usó: una recta \(\ell=(a,b,c)\) es una forma lineal, un punto \(p\) es un vector, y
\(\ell\cdot p=0\) es la incidencia de la sección~\ref{sec:interno}.

\paragraph{La matriz como pila de formas.} Una transformación lineal del plano al plano
devuelve dos números por cada vector; cada uno de ellos es una forma lineal. Con
\[
  A=\begin{pmatrix}2&3\\1&-1\end{pmatrix},\qquad
  A\begin{pmatrix}1\\2\end{pmatrix}=\begin{pmatrix}2\cdot1+3\cdot2\\1\cdot1-1\cdot2\end{pmatrix}
  =\begin{pmatrix}8\\-1\end{pmatrix},
\]
la primera fila mide con la regla \(2x+3y\) y la segunda con la regla \(x-y\).
Geométricamente, aplicar \(A\) es leer dos reglas graduadas a la vez. El núcleo de \(A\)
es el conjunto de puntos donde las dos reglas marcan cero: la intersección de las rectas
\(2x+3y=0\) y \(x-y=0\), que es sólo el origen porque las rectas se cortan
(\(\det A=-5\neq0\)). Si en cambio \(B=\begin{pmatrix}2&3\\4&6\end{pmatrix}\), las dos
filas son la misma regla escalada, sus rectas de nivel cero coinciden, el núcleo es toda
la recta \(2x+3y=0\) y \(\det B=0\): es la sección~\ref{sec:nucleo} vista desde las
filas. Y el método del producto cruz para autovectores (sección~\ref{sec:autovalores})
busca justamente el vector que todas las reglas de \(A-\lambda I\) marcan en cero.

\paragraph{Cómo se transforma una forma cuando se transforma el vector.} Sea
\(A=\begin{pmatrix}1&1\\0&2\end{pmatrix}\) y \(v=(1,1)\). Entonces \(Av=(2,2)\) y
\(\varphi(Av)=2\cdot2+3\cdot2=10\). ¿Qué regla, aplicada al \(v\) original, da ese mismo
\(10\)? Se calcula la fila por la matriz:
\[
  (2,3)\,A=\bigl(2\cdot1+3\cdot0,\ 2\cdot1+3\cdot2\bigr)=(2,8),
  \qquad (2,8)\cdot(1,1)=10\ \checkmark
\]
El vector se multiplica por \(A\) a la izquierda (columna); la forma se multiplica por
\(A\) a la derecha (fila). Es lo que el profesor resume como \(\varphi(A\vec v)=(\varphi A)(\vec v)\):
``una forma lineal se transforma del lado contrario que un vector''.

Si lo que se quiere es \emph{arrastrar la regla junto con la figura}, es decir, encontrar
la forma \(\psi\) cuyas rectas de nivel son las imágenes por \(A\) de las rectas de nivel
de \(\varphi\), hay que deshacer \(A\): \(\psi=\varphi A^{-1}\). Con
\(A^{-1}=\begin{pmatrix}1&-\tfrac12\\0&\tfrac12\end{pmatrix}\),
\[
  \psi=(2,3)\,A^{-1}=\bigl(2,\ -1+\tfrac32\bigr)=\bigl(2,\tfrac12\bigr).
\]
Verificación: el punto \((3,-2)\) está en la recta \(\varphi=0\); su imagen es
\(A(3,-2)=(1,-4)\), y \(\psi(1,-4)=2-2=0\): sigue en la recta de nivel cero de la regla
arrastrada. El punto \((3,0)\), con \(\varphi=6\), va a \((3,0)\) y \(\psi(3,0)=6\).
Hecho geométrico: al deformar el plano las rectas de nivel se deforman con él, siguen
siendo rectas paralelas equiespaciadas (es una afinidad), pero la separación y la
inclinación cambian, y por eso los coeficientes de la forma cambian ``al revés'' que los
de los vectores.

\paragraph{El gradiente es una forma lineal.} Sea \(f(x,y)=x^2+y\) en el punto \((1,1)\),
donde \(f=2\). Para un desplazamiento chico \((h,k)\),
\[
  f(1+h,1+k)=2+2h+k+h^2\approx 2+\underbrace{(2h+k)}_{\text{forma lineal}} .
\]
La derivada en \((1,1)\) es la regla \(2h+k\): la mejor aproximación lineal de \(f\) cerca
del punto. Su recta de nivel cero, \(2(x-1)+(y-1)=0\), es la tangente a la curva de nivel
\(x^2+y=2\) en \((1,1)\). Que ``el gradiente es perpendicular a la curva de nivel'' es la
misma figura de arriba: el vector \((2,1)\) es perpendicular a la familia de rectas de la
forma. Ir a lo chico (la forma lineal en un punto) para volver a lo grande (la superficie
completa) es lo que hace la tercera lectura con el plano tangente y la métrica.

\begin{resultado}[Dualidad en pocas palabras]
\begin{itemize}
  \item Un vector es una columna; una forma lineal es una fila. Fila por columna es un
  número: la medición.
  \item Hecho geométrico: vector = flecha; forma = familia de rectas paralelas
  equiespaciadas (una regla). En \(\R^3\), planos paralelos.
  \item Una matriz es una pila de formas (sus filas) o una hilera de vectores (sus
  columnas). Ambas lecturas son útiles: las filas dan el núcleo, las columnas dan la imagen.
  \item Al transformar el plano con \(A\), los vectores van con \(A\) y las formas con
  \(A^{-1}\) (o con \(A\) del otro lado, si se pregunta ``qué regla mide lo mismo'').
  \item En homogéneas, los puntos son vectores y las rectas son formas; cambiarlos de rol
  es la dualidad de la geometría proyectiva.
\end{itemize}
\end{resultado}

